% Einordnung in den aktuellen Forschungsstand
\section{Einordnung in den aktuellen Forschungsstand}

Die Dash\-board"=Lösung ist eine nicht-invasive, risikoarme Modernisierungsmethode. Es hilft Anwendern, Monitoring-Prozesse effizienter zu gestalten, da es die alten Strukturen in eine nutzerfreundliche Oberfläche verpackt.
Weiterhin ermöglicht das Dash\-board durch dessen informativen Charakter einen vereinfachten Einstieg und unterstützt somit die Einarbeitung neuer Mitarbeiter, was den Generationenwechsel zusätzlich fördert.
%Evtl könnte man noch sagen, dass die Einarbeitung neuer  Mitarbeiter und somit auch der Generationenwechsel damit einfacher geht/gehen  könnte, je nach Faktenlage
Da diese Dash\-board"=Lösung jedoch keine bestehenden Strukturen entfernt, entsteht auch keine Hürde für langjährige Mitarbeiter, welche potenziell Probleme mit der Umstellung haben.

Weiterhin bietet es sich an, eine Ergebnisanzeige mit dieser Lösung zu modernisieren, um vergessene oder geduldete Altlasten aufzudecken. Da es sich teilweise schwierig gestaltet, diese durch die Modernisierungsmaßnahme automatisiert auszulesen, bietet der Prozess eine geeignete Gelegenheit, die Altlasten neu zu entdecken und aufzuarbeiten.

Es kann darüber hinaus dazu dienen, auf das Thema der Modernisierung aufmerksam zu machen, wie diese beispielsweise zur Optimierung von Prozessen eingesetzt werden kann. Das Dash\-board kann somit als eine geeignete und zukunftsfähige Modernisierungsmethode für Legacy-Anwendungen mit Er\-geb\-nis\-an\-zeige eingeordnet werden, welche das Potenzial hat, Altsysteme nachhaltig zu verbessern.

Aufgrund dieser Ergebnisse ist somit bewiesen, dass die Einführung dieser Dash\-board"=Lösung sinnvoll ist, um Legacy"=Monitoring"=Prozesse effizienter und nutzerfreundlicher zu gestalten.

% Feedback:
% Finde gut, dass du Generationenwechsel und Einarbeitung Fachfremder Personen hinzugefügt hast. In der Einleitung stand noch was von Nachhaltigkeit  und Zukunftsfähigkeit. Optional könnte man diese Stichworte hier nochmal aufgreifen, um einen besseren Bezug zur Einleitung zu erhalten.
% Andererseits könnte es sich dann wie eine Wiederholung anfühlen. Dieser Vorschlag hat daher unterste Prio.
% Als Referenz der  Satz aus der Einleitung 1.1: "Diese Forschungsarbeit ist darin motiviert, einen Beitrag zu nachhaltigen und zukunftsfähigen Lösungen beizusteuern."